\section{Logarithmic Functions}
\subsection*{EXPONENTIAL FUNCTION 1 --- B3-10.0}
\addcontentsline{toc}{subsection}{EXPONENTIAL FUNCTION 1 --- B3-10.0}
\label{sec:exponential_function_1}

\subsubsection*{PURPOSE}
The program \texttt{B3-10.0} calculates the exponential function $K^x$, where K = 2, $e$, or 10 and $x: -1 \le x \le 1$. \texttt{B3-10.0} expects $x$ in the accumulator, and returns $K$ in the accumulator.

\subsubsection*{STORAGE REQUIREMENTS}
\begin{center}
\begin{tabular}{ll}
    \toprule
    Memory & 63 sectors (63 words) \\
    Temporary storage & None \\
    \bottomrule
\end{tabular}
\end{center}

\subsubsection*{INPUT}
The input value $x$ must be scaled to $-1 \le x \le 1$ and have a $q$ at $q=1$.

\subsubsection*{CALLING SEQUENCE}
Each of the exponential functions has a separate entry point into the subroutine:

\paragraph{$2^x$}
\begin{table}[h!]
\centering
\begin{tabularx}{\textwidth}{cccX}
\toprule
\textbf{Memory Location} & \textbf{Instruction} & \textbf{Address} & \textbf{Meaning} \\
\midrule
$\alpha-1$ & \texttt{B} & $N$ & Load $N$ at $q=1$ into the accumulator. \\
$\alpha$ & \texttt{R} & \texttt{L+9} & Set return address. \\
$\alpha+1$ & \texttt{U} & \texttt{L} & Call the $2^x$ routine. \\
\bottomrule
\end{tabularx}
\end{table}

\paragraph{$e^x$}
\begin{table}[h!]
\centering
\begin{tabularx}{\textwidth}{cccX}
\toprule
\textbf{Memory Location} & \textbf{Instruction} & \textbf{Address} & \textbf{Meaning} \\
\midrule
$\alpha-1$ & \texttt{B} & $N$ & Load $N$ at $q=1$ into the accumulator. \\
$\alpha$ & \texttt{R} & \texttt{L+9} & Set return address. \\
$\alpha+1$ & \texttt{U} & \texttt{L+2} & Call the $e^x$ routine. \\
\bottomrule
\end{tabularx}
\end{table}

\paragraph{$10^x$}
\begin{table}[h!]
\centering
\begin{tabularx}{\textwidth}{cccX}
\toprule
\textbf{Memory Location} & \textbf{Instruction} & \textbf{Address} & \textbf{Meaning} \\
\midrule
$\alpha-1$ & \texttt{B} & $N$ & Load $N$ at $q=1$ into the accumulator. \\
$\alpha$ & \texttt{R} & \texttt{L+9} & Set return address. \\
$\alpha+1$ & \texttt{U} & \texttt{L+3} & Call the $10^x$ routine. \\
\bottomrule
\end{tabularx}
\end{table}
Each of the entry points is a standard LGP-21 subroutine call; mainline execution resumes directly following the unconditional jump to the subroutine.

\subsubsection*{OUTPUT}
On return, the function value is in the accumulator at $q=4$.

\subsubsection*{REMARKS}
\begin{itemize}
    \item \textbf{Accuracy:} No error is greater than 5 x $10^{-8}$.
    \item \textbf{Time required:} approx. 775 ms.
    \item \textbf{Program format -} the tape contains the program in two versions:
    \begin{itemize}
      \item Relocatable hexadecimal, valid for J1-10.0 (section \ref{sec:PIRloader})
        \item Relocatable decimal, in coding sheet format.
    \end{itemize}
    \item \textbf{Required input/output devices:} None.
    \item \textbf{Overflow:} The program neither checks overflow upon input nor resets it.
\end{itemize}

\subsubsection*{Translator's Notes}
While the input to these functions is strictly restricted to $q=1$, the routine returns the final calculated value in the accumulator scaled to \emph{at $q=4$}. This choice of scale factor represents a best-case compromise with the LGP-21's fixed-point architecture.

Because the maximum possible input value to any of these functions is 1.0, the maximum possible outputs are predetermined:
\begin{align*}
    2^1  &= 2.0 \\
    e^1  &= 2.718281828\dots \\
    10^1 &= 10.0
\end{align*}
Recalling that the representation of the largest possible value for any of these functions as an integer is 10, or $1010_2$ --- four bits in binary --- this means we need to have four bits in front of the binary point, but not more, to represent the maximum value of 10.0: a $q$ of $q=4$.

If we instead chose to return the data at a standard library scale like at $q=9$, this would throw away bits on the right side of the binary point to provide space to on the left-hand side of the binary point to represent values much larger that we would ever need. If, on the other hand, we chose to use $q=1$, this would limit the return value's maximum to 1.99, causing an overflow on $2^1$ and most positive $e^x$ and $10^x$ calculations.

By standardizing the output to at $q=4$, \texttt{B3-10.0} allows the maximum possible value (10.0) for any of these functions to be returned while preserving an optimal 26 bits of precision for the fractional portion of the result for smaller numbers.

\newpage
\subsection*{LOG 1 --- B3-11.0}
\addcontentsline{toc}{subsection}{LOG 1 --- B3-11.0}
\label{sec:logarithm_function}

\subsubsection*{PURPOSE}
The program \texttt{B3-11.0} calculates the logarithm of a positive number to the base 2, $e$, or 10. A 7th-degree polynomial is utilized as the approximation method.

\texttt{B3-11.0} uses an altered multi-argument LGP-21 function call sequence. The function argument must reside in the accumulator, but the two words immediately following the jump to the subroutine must contain the $q$ of the argument and the desired base of the logarithm. The return value
is in the accumulator on return.

\subsubsection*{STORAGE REQUIREMENTS}
\textbf{Program Space:} 1 track, 58 words \\
\textbf{Temporary Storage:} None

\subsubsection*{INPUT}
The required input consists of:
\begin{enumerate}
    \item $N$: The positive number whose logarithm is to be calculated.
    \item $Q$: The $q$ factor of $N$.
    \item $K$: The desired base of the logarithm.
\end{enumerate}

$N$ must be a positive number greater than zero ($N > 0$) and must reside in the accumulator with a positive $q$ factor when branching to the subroutine.

The $q$ factor must lie within the range $0 \le q \le 31$. In the calling sequence, $Q$ is stored in the sector portion of a \texttt{Z} instruction that immediately follows the branch instruction.

$K$, the desired base of the logarithm, is also stored in the sector portion of a \texttt{Z} instruction.
\begin{itemize}
    \item For $\log_2 N$, $K = \texttt{00}$
    \item For $\log_e N$, $K = \texttt{01}$
    \item For $\log_{10} N$, $K$ can take any value other than \texttt{00} or \texttt{01}. 
\end{itemize}
The \texttt{Z} instruction containing $K$ is the final entry in the calling sequence.

\subsubsection*{CALLING SEQUENCE}
\begin{table}[h!]
\centering
\begin{tabularx}{\textwidth}{cccX}
\toprule
\textbf{Memory Location} & \textbf{Instruction} & \textbf{Address} & \textbf{Meaning} \\
\midrule
$\alpha-1$ & \texttt{B} & $N$ & Load $N$ at $q=1$ into the accumulator. \\
$\alpha$ & \texttt{R} & \texttt{L+24} & Set return address. \\
$\alpha+1$ & \texttt{U} & \texttt{L} & Call the log routine. The following parameter constants are effectively scaled at $q=29$. \\
$\alpha+2$ & \texttt{Z} & \texttt{00}$Q$ & The $q$ value of $N$. \\
$\alpha+3$ & \texttt{Z} & \texttt{00}$K$ & Desired base $K$. \\
\bottomrule
\end{tabularx}
\end{table}
As usual, any instruction sequence leaving $N$ in the accumulator prior to the call is fine as long as it leaves $N$ there with a positive $q$. The instruction at location $\alpha$ stores the address ($\alpha+2$) of the \texttt{Z} instruction containing $Q$ into the subroutine's internal tracking mechanism, giving the subroutine the address of $Q$ and $K$.\footnote{Translator's note: the subroutine figures out the return address by taking the address supplied by the \texttt{R} instruction and adding 2 to skip over the two parameter words. It then updates a \texttt{U} instruction in \emph{itself} to jump back to the proper return address}.

Execution resumes following the second parameter constant upon return from the subroutine.

\subsubsection*{OUTPUT}
Upon the return branch to the main program, the calculated logarithm to the desired base resides in the accumulator at $q=6$.

\subsubsection*{OPERATING DETAILS}

\begin{enumerate}[leftmargin=1.5em, itemsep=6pt]
    \item \textbf{Accuracy:} \\
    Any potential error is less than $3 \times 10^{-8}$.

    \item \textbf{Execution time:} \\
    Approximately 1350 ms.

    \item \textbf{Program input:} \\
    The program tape contains the routine in two formats:
    \begin{itemize}
      \item Relocatable hexadecimal, valid for J1-10.0 (section \ref{sec:PIRloader})
        \item Relocatable decimal, in coding sheet format.
    \end{itemize}

    \item \textbf{Required input/output devices:} \\
    None.

    \item \textbf{Overflow:} \\
    Overflow is neither checked upon entry to the subroutine nor reset upon exit.

    \item \textbf{Error Stops:} \\
  \begin{tabularx}{\textwidth}{lX} % Adjust the 9cm width to fit your page margins
    \toprule
    \textbf{Memory Location} & \textbf{Meaning and Remedy} \\
    \midrule
    \texttt{main program+8}  & Argument $N$ is zero or negative. The accumulator contains the value $N - (1$q=30$)$. To continue, deposit a corrected value, and jump manually via the console to resume execution. \\
    \bottomrule
  \end{tabularx}
\end{enumerate}

\subsubsection*{Translator's Notes}
The behavior on an invalid argument ($N \le 0$) is a little odd. Halting at \texttt{L+8} leaves execution close to the error, fine; but why also leave $N - (1\text{ at }$q=30$)$ in the accumulator instead of just $N$?

``$1$ at $q=30$'' means a single \texttt{1} bit in position 30 of the word, i.e.\ the value $2^{-30}$ --- one bit short of the LSB. Subtracting it from $N$ is essentially $N - \epsilon$: the result looks almost identical to $N$ if you read it as a number, but the low end of the word will differ.

The LGP-21's console isn't a wall of indicator lights; it's a small oscilloscope tube that paints the accumulator, the address register, and the current instruction as rows of dashes --- a \texttt{0} bit is drawn low, a \texttt{1} bit drawn high. The machine is slow enough that on a stop, the row sits there steady and you can read off the bits.

So the actual point of \texttt{L+8} probably isn't that the difference makes the value easier to recognize visually --- it's that subtracting that one bit is a cheap way to \emph{prove}, at the stop, that the routine got as far as detecting the error and ran the trap. If the operator sees a value with a flipped low end, the error path ran; if they see the original $N$ untouched, something else stopped the machine. The accumulator doubles as a one-bit status flag without spending a memory cell on one.
\newpage
